\section{Finite-degree completion}\label{sec:finite}
For real $a$ and positive $t$, write $\mathcal A(a,t)$ for the assertion
that every finite simple graph with $n$ vertices and $m$ edges has a packing
of weight at most $an+m/t$. This section supplies the finite-degree cases
of the main induction and the completion of the stopped remainder.

\subsection{The short-circumference starting point}
\begin{lemma}[Short-circumference certificate]\label{lem:six}
If $G$ has no simple cycle of length at least $7$, then it has a packing
$\mathcal P$ satisfying
\[
                  z(\mathcal P)\le19r(G)+6m.
\]
\end{lemma}
The proof is a structural reduction followed by a finite certificate check;
it is given in Appendix~\ref{app:terminal}. Both the reduction to the finite
catalogue and the interpretation of its witnesses as actual original-edge
packings are part of the formal proof.

\begin{lemma}\label{lem:harmonic}
For every integer $u\ge7$,
\begin{equation}\label{eq:harmonic}
 \mathcal A\left(\frac{19}{42}+\frac12\sum_{j=8}^u\frac1j,u\right)
 \quad\hbox{holds}.
\end{equation}
\end{lemma}
\begin{proof}
Greedily delete simple cycles of length at least $7$. Their number is at
most one seventh of the number of deleted edges. On the remaining graph,
Lemma~\ref{lem:six}, divided by $42$, gives weight at most
$19n/42+m'/7$. Adding the deleted cycles proves $\mathcal A(19/42,7)$.
For the successor step from $t$ to $t+1$, first remove cycles of length at
least $t+1$. On the remaining graph, the Erd\H{o}s--Gallai long-cycle edge
bound~\cite{EG1959} gives $2m'\le tn$. Consequently
\[
 \frac{m'}t-\frac{m'}{t+1}
          =\frac{m'}{t(t+1)}\le\frac{n}{2(t+1)}.
\]
Apply $\mathcal A(a,t)$ to this remainder and pay the deleted cycles by
their lengths at denominator $t+1$. This proves
$\mathcal A(a+1/(2(t+1)),t+1)$.
Iterating that step from $7$ to $u$ gives~\eqref{eq:harmonic}.
\end{proof}

\subsection{The mixed cover}
Put $B_*=43/10$, and define
\[
 D(n,m)=B_*nm^{1/3}+3n^{3/2},\qquad
 \psi(v)=\begin{cases}(53/200)v,&0\le v\le1,\\
              1/2-(47/200)v^{-1/3},&v\ge1.
              \end{cases}
\]
The two formulas for $\psi$ agree at $1$. The function is nonnegative,
nondecreasing and concave, and is at most $1/2$.

\begin{lemma}[Mixed-cover certificate]\label{lem:mixed}
Every graph of positive order has each of the following packings:
\begin{align}
 p&\le n\psi\left(\frac{m}{2B_*^3n}\right),
       &q&\le D(n,m)+2n;\label{eq:fine-cover}\\
 p&\le n/2,&q&\le D(n,m)+n.\label{eq:coarse-cover}
\end{align}
The packings in the two lines need not be the same.
\end{lemma}
\begin{proof}
Use strong induction on the order. If $m\le D(n,m)$, the empty packing
works. The greedy range is handled by the long-cycle estimate. A sparse
cut is paid from the difference between the parent budget $D(n,m)$ and
the two child budgets; cycles in the two induced sides are assembled and
the cut edges are left as singletons. The perspective of the concave
function $\psi$ pays the cycle counts in~\eqref{eq:fine-cover}.
A sufficiently low-degree vertex is handled by deletion and the corresponding
one-vertex budget inequality. The remaining graphs meet the degree and cut
hypotheses of the routed-sampling construction. Its certificate covers
\emph{all} the remaining continuous parameters, including the unbounded
scale range. Appendix~\ref{app:cover} gives the construction, the precise
coverage interface, and the finite checks closing this induction.
\end{proof}

\subsection{A certified depth-first-search step}
Consider a depth-first-search block with $s$ new vertices and at most $T$
boundary vertices, and let $m'$ denote its assigned edges. The boundary is
padded to $T$ vertices when needed. Choose the local packings first, and
then apply $\mathcal A(a,t)$ once to the graph of all uncovered original
edges, retaining the original vertex set. Its intercept is therefore $an$.
For a regular block, the fine packing contributes the following additional
fee, after subtracting the block's assigned-edge budget $m'/T$:
\begin{equation}\label{eq:regular-fee}
 (s+T)\psi\left(\frac{m'}{2B_*^3(s+T)}\right)
       +\frac{D(s+T,m')+2(s+T)}t-\frac{m'}T.
\end{equation}
The alternative is to leave the block's edges for that common suffix
estimate; its additional fee is $m'/t-m'/T$. An entire component with no boundary instead uses
\eqref{eq:coarse-cover}, giving fee
\begin{equation}\label{eq:component-fee}
                  s/2+\frac{D(s,m')+s}t-\frac{m'}T.
\end{equation}
These are different branches of the DFS decomposition, not interchangeable
versions of the same packing.

For $g>0$ let $\Pi_p(g)=(1+g)^pg^{1-p}$. Write
\begin{align}
 \mathcal K(T,t,d,g)
  ={}&(1+g)\left[\psi\left(\frac{d}{4B_*^3}\right)
                              +\frac2t-\frac d{2T}\right]\notag\\
     &+\frac{B_*}{t}\left(\frac{dT}{2}\right)^{1/3}\Pi_{4/3}(g)
                +\frac{3\sqrt T}{t}\Pi_{3/2}(g).\label{eq:core-fee}
\end{align}
The substitution $d=2m'/(s+T)$ and $g=T/s$ makes
\eqref{eq:regular-fee} equal to $s\mathcal K(T,t,d,g)$.
The functions $\Pi_p$ are convex for $p\ge1$: they are perspectives of
$x\mapsto x^p$. Thus, for fixed $d$, the maximum of~\eqref{eq:core-fee}
on a closed $g$-interval occurs at an endpoint, even when its linear
coefficient is negative.

\begin{lemma}[Finite DFS step]\label{lem:dfs-step}
Let $T,t$ be positive integers with $T\ge2^{20}$, and let $\beta_0,K>0$.
Suppose
\begin{align*}
 &(1+1/\beta_0)/2\le K,\\
 &\mathcal K(T,t,d,g)\le K\quad(d>t),\quad
       g\in\{(2\beta_0+2^{-20})^{-1},\beta_0^{-1}\},\\
 &s/2+\bigl(D(s,m')+s\bigr)/t-m'/T\le Ks
                 \quad(0<s\le\beta_0T+1,\ m'\ge0).
\end{align*}
Then $\mathcal A(a,t)$ implies $\mathcal A(a+K,T)$.
\end{lemma}
\begin{proof}
First remove simple cycles of length at least $T$. If $h$ edges are
removed, these cycles cost at most $h/T$. The remaining graph, on the same
vertex set, has $m-h$ edges and no cycle of length at least $T$; this is the
circumference hypothesis required by the DFS block construction.
The rounded DFS construction assigns each remaining edge to exactly one block.
Its regular blocks have $T/s$ between the two displayed endpoints; the
rounding error is bounded by $1/T\le2^{-20}$. For $d\le t$, the direct
suffix branch is paid by $(1+1/\beta_0)/2\le K$. For $d>t$, convexity
and the two endpoint hypotheses pay~\eqref{eq:regular-fee}. Boundary-free
components have at most $\beta_0T+1$ vertices and are paid by the last
hypothesis. Sum over the disjoint sets of new vertices. The assigned edge
counts sum to $m-h$. Apply the suffix packing once to the actual residual
graph, so its intercept is charged once per original vertex. This gives
weight at most $(a+K)n+(m-h)/T$ on the remainder. Restoring the initially
removed cycles adds at most $h/T$, yielding $(a+K)n+m/T$.
The original-edge assembly is proved by
\mcode{Cycle13.weightedAffine_dfs_step}; the scalar hypotheses are supplied in
\mcode{Cycle12.FiniteComponentPayment}.
\end{proof}

The endpoint hypotheses in Lemma~\ref{lem:dfs-step} are verified on whole
real density intervals. In the high-density branch, $t\ge4B_*^3=79507/250$.
After substitution, the nonlinear density expression has the form
$A x^{1/3}-B x^{-1/3}-x/2$; a certified tangent and derivative sign pay its
entire domain. Rational lower and upper bounds for $A,B$ and $\sqrt{\cdot}$
are established by cubing or squaring nonnegative rational endpoints.
The component branch has $B=0$ and also includes zero density.
No finite sampling of possible graph densities is used.

There are $83$ finite chains. A chain begins with~\eqref{eq:harmonic}, uses
a certified upper bound on its harmonic intercept, and applies
Lemma~\ref{lem:dfs-step} a finite number of times. If its final affine bound
is $\mathcal A(a,T)$ and its degree cap is $d_0$, then
\begin{equation}\label{eq:finite-completion}
       2m\le d_0n\quad\Longrightarrow\quad
       \wt_\sigma(\mathcal P)\le\left(a+\frac{d_0}{2T}\right)n.
\end{equation}
All $T$ and $d_0$ in the formal certificates are integers; bounds such as
$2^x\le d_0$ are certified before~\eqref{eq:finite-completion} is applied.
The chains and their exact fees are tabulated in the companion data.

\begin{proposition}[Finite potential base]\label{prop:finite-base}
If $2m\le2^{B_{\rm main}}n$, then $G$ has a packing of weight at most
$C\phi(n)n$.
\end{proposition}
\begin{proof}
The $45$ base intervals cover the required order range, with their shared
endpoints included. On each interval choose its certified finite chain.
Use the simple-graph estimate $2m\le n^2$ together with the assumed cap
to establish the chain's integer degree cap. The resulting constant
in~\eqref{eq:finite-completion} is at most the lower endpoint value of
$C\phi(n)$. Monotonicity from Lemma~\ref{lem:potential} pays the entire
interval. Beyond the last interval the main finite chain and the same
monotonicity apply. The zero-order case has no edges. The endpoint payments
and the interval-coverage proof are
\mcode{Cycle12.FinitePotentialPayments} and
\mcode{Cycle12.FinitePotentialBase12}.
\end{proof}
